
Uncertainty quantification (UQ) methods have been proposed for automatic hallucination detection in large language models, yet prior evaluations employed heterogeneous conditions that obscure method-specific performance characteristics. This work presents a controlled head-to-head comparison of token-level and semantic-level UQ methods on TruthfulQA multiple-choice hallucination detection using Llama-3-8B-Instruct. Contrary to expectations derived from prior work, semantic entropy---which clusters semantically equivalent responses via natural language inference before computing entropy---achieves an AUROC of 0.5645, while simple max probability achieves 0.8068. Mechanism verification confirms that the NLI-based clustering functions correctly (average 4.92 clusters from 5 samples, entropy standard deviation 0.0752). The 0.24-point AUROC gap stems from format mismatch: multiple-choice answers consist of single letters (A/B/C/D) that provide insufficient semantic content for meaningful NLI comparison. Token-level methods extract discriminative signals directly from logits regardless of answer format. These findings indicate that UQ method effectiveness is format-dependent, and practitioners should select methods appropriate to their task format rather than assuming more sophisticated approaches yield better performance.

